\section{Evaluation}
\label{sec:evaluation}

\noindent\textbf{Evidence status.} The following four packages are reserved
experiment designs and result slots. They specify the minimum evidence needed
for the paper's claims; they do not report completed experiments. The historical data were previously used in a train/validation split. We
therefore use a retrospective chronological benchmark, disclosing prior use
and fixing this revision's train/validation/test intervals before execution.
Fit on training and select on validation; do not revise methods using the
new test outcomes. This protocol does not create a previously unseen holdout.

\subsection{E1: Are the Work and Contention Scenarios Sound?}
\slot{E1: inputs and scenario contract}{
Use Figure~12 of AMSP's March 2024 revision~\cite{chen2024amsp}: LLaMA-7B/13B/30B,
BF16, 4096-token sequences, and fixed global batch. Its 32/128/512/1024-GPU
points map to 4/16/64/128 eight-A800 nodes. Archive PDF hash, vector-bar
extraction, units, and the simulation configuration. Audit input workloads,
state initialization, work accounting, and independent wait probes.
Scenario statistics and validation outcomes: \pending{E1}.}

\paragraph{What the traces represent.}
A background job supplies a tuple of submission time, node demand, requested
wall time, and occupied duration. These tuples define rigid, homogeneous-node
workloads with no modeled network or I/O interference. Each source's time
order and joint width--duration structure remain intact; traces are not pooled.
The nominal cluster has 128 nodes, matching the profile testbed's size. This
is a declared experimental capacity, not an estimate of either source cluster.
Old durations specify scenario demands; they do not predict the old programs'
runtimes on A800s. Queue state and waits are regenerated under the declared
scheduler. Conditional on identical tuples, initial state, and scheduler,
changing hardware labels cannot change admission. That conditional property
neither makes workloads hardware-independent nor establishes universal transfer.

\paragraph{Replay and validation.}
The main scheduler is a node-level conservative-backfill approximation, not
Slurm itself. Rebuild background state from an empty state at the declared
trace-prefix origin, excluding at least 28 days before scored arrivals and
checking prefix sensitivity in development. Do not transplant historical starts
or running allocations into the constructed cluster. Methods branch from the
same background state; their requests then affect subsequent admissions.
State continues across all chunks. Original observed waits are not ground
truth for this constructed scenario. Verify scheduling mechanics on
independently solvable cases and evaluate the predictor on held-out simulator
probes; source-calendar tests preserve chronology and report censoring.

\begin{table}[t]
\centering\small
\begin{tabular}{@{}p{0.42\linewidth}p{0.50\linewidth}@{}}
\toprule
Input / check & Reserved evidence \\
\midrule
Trace demand and burstiness & Pending: counts, widths, durations \\
Scenario contention & Pending: occupancy and wait tails \\
Published throughput input & AMSP Fig.\ 12; extraction archived \\
State and work conservation & Pending: scheduling/accounting checks \\
Prediction vs.\ probe wait & Pending: errors and coverage \\
Unmeasured execution costs & Declared overhead scenarios \\
\bottomrule
\end{tabular}
\caption{Input and validity results placeholder. Report both source workloads
separately, including time zones, cleaning, clipping, and calendar blocks.
Probe labels come from the simulator; they are not observed alternative
requests on a production cluster.}
\label{tab:validity}
\end{table}

\paragraph{Execution inputs.}
Figure~12 reports tokens per GPU per second, $z_n$; convert to updates per
hour as $q_n=3600(8n)z_n/(1024\cdot4096)$. The 1024-sequence batch follows
from the source's stated micro-batch and accumulation settings. Extract only
published points; 64- and 256-GPU points are absent. For each model, freeze
$U_\star=\lceil192q_4\rceil$ before scheduling experiments. This chooses a
workload of approximately 192 useful-compute hours on four nodes; overhead
and rounding determine the chunk count, rather than a convergence target. Use 600 seconds of total allocated overhead per chunk, split equally
between initialization/restart and saving, as an explicit scenario. E4 varies
this unmeasured cost. No checkpoint implementation, final-quality equivalence,
or complete-node power measurement is claimed.

\subsection{E2: Is Adaptation Better than Strong Alternatives?}
\slot{E2: primary comparison}{
Evaluate paired initial arrivals for both source workloads and all three published profiles.
Use EIA's Texas ERCOT average consumption-based CO$_2$ intensity as the primary
scenario~\cite{eia2026gridmonitor}. Set $T_{\rm ref}$ to the fastest fixed
training mean TAT and $D_{\rm hi}$ to the training Fixed-4 95th-percentile TAT.
Use $\beta_{\min}=1$, $\beta_{\max}=D_{\rm hi}/T_{\rm ref}$, and
$s\in\{0,1/4,1/2,1\}$. This sets the budget span, not a feasibility guarantee.
Freeze the grid, miss tolerance, power interval, seeds, and selection rule before test access. Compare \name{} with every fixed scale, a validation-fitted
fixed mixture, plan-once control, and rollout planning. Main figure, carbon ratios,
miss rates, TAT, and node-hours: \pending{E2}.}

\paragraph{Baselines.}
Fixed-4/16/64/128 retain the same scale across all chunks and use the identical
work, checkpoint, and final-request rules. \emph{Best-Fixed} is selected on
validation for each budget. \emph{Fixed-Mix} samples one scale at the start of
a run and keeps it: its four mixture weights solve the same endpoint objective
and miss constraint using validation statistics. Mixture feasibility depends
on the mixture of miss probabilities, not on an interpolated mean TAT.
All four fixed points remain visible even if dominated.

For the planning baselines only, fit a request-conditioned boosted regressor
for $\log(1+w)$; chronological training residuals supply 32 equal-mass atoms
of $\widehat F_{k,n}$. Its causal exposure estimate is
\begin{equation}
 \widehat L_{k,n}=n\,\mathbb E_{w\sim\widehat F_{k,n}}
 \left[\int_{t_k+w}^{t_k+w+d_{k,n}}\widehat c_{t_k}(s)\,ds\right].
 \label{eq:forecast}
\end{equation}
E1 reports probe errors and coverage to document this approximation; prediction
accuracy is not an acceptance criterion for \name{}.

\emph{Rollout-MPC} enumerates two upcoming scale choices and a constant-scale
tail to completion, uses the auxiliary wait model and causal CI forecast, and executes
only the first action before replanning. There are at most $4^3$ candidate
plans. Future queue observations are unavailable, so the rollout holds the
current queue and wait-model calendar features fixed, updates request length
and the forecast integration window, and samples conditional waits. Its carbon
estimate adds past endpoint costs using released CI, retaining each chunk's
submission forecast for unreleased intervals, to the planned continuation. It selects the
lowest worst-endpoint normalized carbon plan meeting the estimated miss tolerance, or the
plan with the smallest miss estimate if none qualifies. This is a strong
causal baseline with an explicit approximation, not a clairvoyant scheduler.
\emph{Plan-once} runs the identical planner only at the first submission and
executes its selected scale sequence and constant tail without replanning.
It receives the initial queue and forecast, allowing a planned within-run
mix while removing feedback from later admissions. Both planning variants
may tune their internal risk threshold on validation; test compares them
under the same actual miss tolerance as every other method.

\begin{figure*}[t]
\centering
\figslot{1.65in}{E2: Two source workloads $\times$ three published AMSP profiles.\\
All fixed scales; Fixed-Mix; Plan-once; Rollout-MPC; \name{}.\\
Mean TAT versus modeled carbon at a declared $\rho$; mark observed miss violations and label slider ticks.\\
Companion labels: worst-endpoint carbon ratio and deadline-miss rate.}
\caption{Primary results placeholder. Every point completes the same prescribed
work from paired initial arrivals. Report hours, modeled carbon normalized
within each contention--profile panel, 95\% paired block intervals, the power
scenario, and the frozen budget grid. Mark observed miss-target violations
explicitly. Fixed-Mix is an executable per-run randomized policy; connecting
other points does not imply an evaluated intermediate policy. Report every
slider tick; do not smooth reversals or select settings after test.}
\label{fig:main}
\end{figure*}

\paragraph{Paired analysis and selection.}
An episode is one initial arrival and complete run. Resample contiguous calendar
blocks jointly across methods, with RL seeds as another replication level.
Blocks must cover the history window and development dependence. Training fixes
budgets and normalizers; validation fixes operating points. The primary
comparator is the strongest validation-selected non-RL method meeting the
empirical miss target at the same budget. Fixed-128 is not the automatic
denominator. Confidence bounds are reported separately from this frozen
empirical selection rule.

Report the endpoint ratio of \emph{mean} carbon costs, mean and p95 TAT,
miss rate, and allocated node-hours per panel; macro-averages use dimensionless
panel ratios. Carbon includes all completed runs, including deadline misses.
Episodes crossing a split must be purged or assigned entirely to one split,
and any run not observed to completion is reported as censored, not silently
dropped. Texas CI must have a documented source, release time, time zone,
resolution, and missing-data rule. If CI and queue years do not overlap,
their alignment is a labeled counterfactual pairing.

\paragraph{Result slot and inference.}
\pending{E2: effect sizes, intervals, eligible budgets, and failed panels}.
An observed budgeted carbon reduction requires both frozen methods to meet
the miss target on complete validation and test cohorts. This describes
observed tradeoffs, not population feasibility. A stronger statistical claim
also requires joint endpoint cost intervals and one-sided miss bounds to pass
the declared criteria. Report inconclusive bounds, including for zero misses;
nearby arrivals do not create independent calendar blocks. Gains over
Fixed-128 alone establish neither adaptive nor RL value. An inspected former
test period cannot be relabeled as untouched.

\subsection{E3: Does Feedback Beat a Learned Sequence?}
\slot{E3: matched feedback comparison}{
Reuse E2 \name{}, Plan-once, and Rollout-MPC. Train one additional
\emph{Precommitted-RL} policy with the same architecture, budget grid, seeds,
arrival sampling, and episode-interaction budget. Before its first submission,
it samples an entire scale sequence using only initial public queue/CI inputs,
remaining planned work, and $D$ minus planned allocation time. It receives no
later queue observations, CI releases, or realized elapsed time. Evaluate at
one prespecified budget per panel. Results: \pending{E3}.}

\begin{table}[t]
\centering\small
\begin{tabular}{@{}lccc@{}}
\toprule
Method / variant & Carbon & Misses & Node-h \\
\midrule
\name{} & --- & --- & --- \\
Precommitted-RL & --- & --- & --- \\
Rollout-MPC & --- & --- & --- \\
Plan-once & --- & --- & --- \\
\bottomrule
\end{tabular}
\caption{Feedback results placeholder at prespecified budgets. All dashes
are pending E3 results. Use E2 normalization and paired intervals; compare
carbon jointly with misses. Precommitted-RL changes its scale within a run
but commits its choices before seeing subsequent outcomes.}
\label{tab:ablation}
\end{table}

The learned-sequence comparison removes feedback while retaining the capacity
to mix scales and optimize from completed episodes. Its remaining-time input
is planned allocation time only, not a zero-wait assumption in the execution
environment: all actual waits and costs enter training returns and evaluation.
Archive each full plan before executing it. Neither policy needs an admission
predictor. Matching interaction budgets prevents extra replay from explaining
a feedback gain, although it cannot guarantee equal optimization quality.
Plan-once versus MPC separately tests replanning within an explicit planner.
Report single-chunk frequency: those runs offer no later decision. A fixed-scale
win without a learned-sequence win supports mixing, not a feedback advantage;
a win only over MPC could reflect its modeling error. Action frequencies alone
do not establish either claim. Results: \pending{E3: feedback increment}.

\subsection{E4: Does the Conclusion Depend on the Power Model?}
\slot{E4: robustness from existing logs}{
Freeze policies and recompute every E2 execution over the declared $\rho$
interval, including both endpoints. Report pairwise break-even points,
componentwise exposure differences, and the independent scale-error radius
supported by Equation~\ref{eq:box}. Compare the robust policy with a
single-endpoint-trained counterpart at one prespecified budget per panel.
Reuse existing logs for parameter sensitivity; queue and forecast perturbations
require fresh replay. Outcomes: \pending{E4}.}

\begin{figure}[t]
\centering
\figslot{1.05in}{E4: carbon advantage versus $\rho$.\\
Endpoint intervals; zero-gain line; break-even $\rho$.\\
Fixed policies and comparator throughout.\\
Small inset: endpoint-trained versus robust policy.}
\caption{Power-sensitivity results placeholder. The parameter range is an
assumed scenario, not a measured power range. Endpoint comparisons characterize
the interval only under Equation~\ref{eq:power}. Plot uncertainty from paired
execution samples separately from power-scenario variation.}
\label{fig:power}
\end{figure}

Figure~\ref{fig:power} reserves the sensitivity result.
\pending{E4: interval-wide advantage or the exact region of reversal}.
For $\Delta A=A_\pi-A_b$ and $\Delta B=B_\pi-B_b$, a possible break-even point
is $\rho^\star=-\Delta B/(\Delta A-\Delta B)$ when the denominator is nonzero.
A point outside the declared interval is not an in-range reversal.
Simultaneous one-sided endpoint intervals, obtained from the paired block
resamples, support an interval-wide statistical statement; two point estimates
alone do not. A nominal gain that reverses near the selected power assumption
must be reported as model-sensitive.

\paragraph{Does the simulated environment determine the gain?}
At one prespecified budget and the 7B profile, reuse frozen full-policy and
Rollout-MPC models in each source workload. Independently change (i) background
widths to $\min(2n,128)$, reporting the capped fraction, (ii) scheduling to
strict FCFS, and (iii) total chunk overhead to 60 or 1800 seconds. Preserve
physical budgets, workload, CI calendars, and all fitted parameters. These
four additional settings require new replay, not new training or hardware
measurements. Width and scheduler changes test learned queue relationships;
E1 prediction accuracy alone cannot establish this robustness. Report carbon and misses together: a gain that disappears or
violates the budget limits the claim. They test the selected contention and
overhead scenarios, not every possible cluster. Results: \pending{E4: environment sensitivity}.

A compact reuse analysis reports E2 outcomes by held-out calendar block and
CI variability. A constant-CI rescore of the same logs removes timing-related
carbon differences for those executions; it does not simulate policies reacting
to constant CI. No cross-region, multi-user-adoption, or facility-wide
experiment is needed for the scoped claim.
